\section{Requirement Analysis}

\subsection{Literature Survey}

\subsubsection{Theory Associated with the Problem Area}
Multi-object tracking first detects objects and then keeps an ID across frames.
ByteTrack checks high-score boxes and also uses lower-score detections in a
second matching step \cite{zhang2022bytetrack}. This helps during short
occlusion. It does not identify the actual player, so team colour, jersey
evidence, appearance, and manual review are added in our design.

Pitch calibration converts an image point to a field point. The bottom centre
of a player's box is treated as the contact point. A homography matrix $H$ maps
the image point $p=(u,v,1)^T$ to $q$:
\begin{equation}
q \sim Hp, \qquad (x,y)=\left(\frac{q_1}{q_3},\frac{q_2}{q_3}\right).
\end{equation}
Distance, speed, zones, team width, and team length are calculated only after
the mapping error is checked against known field points.

The wearable signals have different rates and errors. The IMU measures rapid
movement, while GPS updates more slowly. PPG measures light changes at the
skin, but football movement can hide the pulse signal \cite{zhang2021ppg}. The
system checks contact, red and infrared signal range, motion, peak consistency,
and sample age before showing BPM or estimated SpO$_2$.

Clock alignment does not use packet arrival order. SNTP anchors the ESP32 UTC
clock to its monotonic timer with
\begin{equation}
t_s = t_d + b,
\end{equation}
where $t_d$ is the monotonic device time, $b$ is the current UTC offset, and
$t_s$ is source time. Source time, relay receive time, local receive time, and
sequence are kept. The remaining alignment error is saved in the joined record.

Team shape uses positions on the field. The centroid gives the team centre,
width and length give the spans on both axes, and compactness gives the spread
around the centre. These measures are common in tactical movement studies
\cite{goncalves2017tactical} and are easy to check on a pitch view.

Workload prediction needs repeated sessions rather than one sensor value.
Inputs can include high-intensity distance, acceleration count, session load,
recovery time, change from personal baseline, and data quality. A prediction
model is tested on later sessions. Its probability quality can be checked with
the Brier score
\begin{equation}
\mathrm{BS}=\frac{1}{N}\sum_{i=1}^{N}(p_i-y_i)^2.
\end{equation}

\subsubsection{Existing Systems and Solutions}
SoccerNet-Tracking is a football benchmark with players, goalkeepers,
referees, the ball, team labels, and jersey data \cite{cioppa2022soccernet}.
SportsMOT adds many examples of fast movement and similar-looking players
across sports \cite{cui2023sportsmot}. They are useful for detector and tracker
work, but they do not contain our wearable data or campus camera view.

Honda et al. combine video and player trajectories for pass-receiver prediction
\cite{honda2022pass}. TacticAI studies player relations in corner kicks
\cite{tacticai2024}. These projects show that tactical context matters, but
their task and data are different from a live campus training platform.

Workload studies use GPS and training data for injury-risk modelling. Rossi et
al. use an interpretable decision-tree approach \cite{rossi2018injury}, while
Pi{\l}ka et al. study XGBoost with wearable load data \cite{pilka2023workload}.
Their reported scores cannot be treated as our score. A review by Leckey et al.
also shows large differences in labels, groups, and testing methods
\cite{leckey2025injury}.

Table~\ref{tab:literature_comparison} compares the main work used in our design.
The contribution column describes the cited work, not a result from this
project.

\begin{longtable}{@{}p{0.95in}p{1.2in}p{1.65in}p{1.75in}@{}}
\caption{Comparison of related work}\label{tab:literature_comparison}\\
\toprule
\textbf{Work} & \textbf{Input} & \textbf{Main contribution} & \textbf{Gap for this project}\\
\midrule
\endfirsthead
\toprule
\textbf{Work} & \textbf{Input} & \textbf{Main contribution} & \textbf{Gap for this project}\\
\midrule
\endhead
ByteTrack \cite{zhang2022bytetrack} & Detections & Matches high- and low-score boxes & No football roster or wearable link\\
SoccerNet-Tracking \cite{cioppa2022soccernet} & Football video & Football MOT dataset and benchmark & No physiological stream\\
SportsMOT \cite{cui2023sportsmot} & Multi-sport video & Sports tracking benchmark & No jersey-device binding\\
Honda et al. \cite{honda2022pass} & Video and trajectories & Pass-receiver prediction & One tactical task, no wearable context\\
TacticAI \cite{tacticai2024} & Corner-kick data & Player-relation prediction and suggestions & Not a full live acquisition system\\
Rossi et al. \cite{rossi2018injury} & GPS training load & Interpretable injury forecasting & Study labels do not transfer directly\\
Pi{\l}ka et al. \cite{pilka2023workload} & Wearable load & XGBoost injury-risk study & No synchronised tactical view\\
Leckey et al. \cite{leckey2025injury} & Published studies & Review of injury-risk methods & Shows need for careful local testing\\
Teixeira et al. \cite{teixeira2025tactical} & Tactical studies & Review of football tactical analysis & Notes data access and explanation gaps\\
\bottomrule
\end{longtable}

\subsubsection{Research Findings from Existing Literature}
The review gives five design rules. Weak detections can help tracking, but they
still need identity checks. Public football data helps early training, but
campus footage needs its own labels. Tactical values must use player relations.
Wearable values need movement-aware quality checks. Workload models must be
tested on later sessions and explained in simple terms.

Our work builds on these findings through one session pipeline. It links a
wearable to a visible player, preserves time and quality during fusion, and
shows the supporting event to the coach. The aim is a testable campus product,
not a replacement for the systems reported in the papers.

\subsubsection{Problem Identified}
The reviewed work does not provide one solution for live player and ball
tracking, jersey support, low-cost wearables, measured time alignment, player
correction, tactical analysis, workload review, and a storage-free relay. The
main problem is correct joining. A good heart-rate reading given to the wrong
player or wrong second is still a wrong result.

\subsubsection{Survey of Tools and Technologies Used}
Python is used for vision, sensor processing, analytics, and the API.
Ultralytics provides the detector interface, ByteTrack provides association,
and OpenCV handles video and pitch mapping. FastAPI and WebSocket provide live
updates. React, TypeScript, and Vite provide the responsive dashboard.

The wearable uses an Espressif ESP32-S3. Arduino CLI and the official
Espressif board package are installed in the \texttt{soccer} conda environment.
The required live relay runs in Docker behind Cloudflare and Caddy at
\texttt{cpg44.nivaspms.com}. PlantUML source is used for diagrams and LaTeX is
used for the report and presentation.

\subsection{Software Requirement Specification}

\subsubsection{Introduction}

\paragraph{Purpose}
The SRS gives requirements that can be checked during the prototype and final
evaluation. It covers capture, synchronisation, analytics, relay transport,
dashboard use, review, and data export.

\paragraph{Intended Audience and Reading Suggestions}
The intended readers are the mentor, panel members, project team, testers, and
future student maintainers. Readers can use this section for requirements and
Section~\ref{sec:design_specifications} for their design.

\paragraph{Project Scope}
The product starts with camera frames and wearable samples. It ends with live
views, reviewed events, session evidence, and training data. Medical diagnosis
and official match decisions are outside the scope.

\subsubsection{Overall Description}

\paragraph{Product Perspective}
The platform is a web application whose main processing and session storage run
on the project PC. The VPS relay is a required deployed component for the public
domain route. It forwards recent authenticated events and keeps no database or
match archive.

\paragraph{Product Features}
Main features are session setup, camera calibration, roster and team setup,
wearable pairing, live player and ball tracking, telemetry quality, tactical
events, workload review, identity correction, replay, and checked export.

\begin{longtable}{@{}p{0.65in}p{4.55in}p{0.65in}@{}}
\caption{Functional requirements}\label{tab:functional_requirements}\\
\toprule
\textbf{ID} & \textbf{Requirement} & \textbf{Priority}\\
\midrule
\endfirsthead
\toprule
\textbf{ID} & \textbf{Requirement} & \textbf{Priority}\\
\midrule
\endhead
FR-01 & Create, start, pause, resume, and close one timestamped session. & Must\\
FR-02 & Calibrate the pitch and report mapping error. & Must\\
FR-03 & Detect and track players, goalkeepers, referees, and the ball. & Must\\
FR-04 & Keep team, jersey, appearance, and manual identity evidence. & Must\\
FR-05 & Read wearable packets with device time, receive time, and sequence. & Must\\
FR-06 & Bind a wearable to a player and keep binding corrections. & Must\\
FR-07 & Align both streams and report the remaining timing error. & Must\\
FR-08 & Show a large live pitch, player data, quality, and stream health. & Must\\
FR-09 & Calculate tactical geometry, movement load, and event summaries. & Must\\
FR-10 & Export reviewed observations, labels, features, and source details. & Must\\
FR-11 & Relay live events only through \texttt{cpg44.nivaspms.com}. & Must\\
FR-12 & Replay a bounded sequence window after a short disconnection. & Must\\
\bottomrule
\end{longtable}

\subsubsection{External Interface Requirements}

\paragraph{User Interfaces}
The dashboard shall support keyboard and pointer use and remain readable on
laptop, projector, and tablet widths. Status shall use text along with colour.
The pitch and live markings shall remain the largest part of the live page.

\paragraph{Hardware Interfaces}
The camera interface accepts a local camera or supported video stream. The
wearable sends timestamped raw samples by HTTPS to
\texttt{cpg44.nivaspms.com}. The local processor receives the samples from the
relay by WSS. Each event includes source identity, time, sequence, payload, and
status.

\paragraph{Software Interfaces}
FastAPI provides session commands and snapshots. Versioned WebSocket events
provide live updates. Remote clients use HTTPS or WSS at the required domain;
direct VPS-IP addresses are rejected.

\subsubsection{Other Non-functional Requirements}

\paragraph{Performance Requirements}
The system shall show processing rate, display delay, packet age, and sequence
gaps for the demonstration setup. Queues and relay memory shall have fixed
limits. A relay reconnect shall request events after the last accepted sequence
or request a new snapshot when the gap is too old.

\paragraph{Safety Requirements}
The wearable shall be insulated, strain-relieved, and fitted so it does not
interfere with play. A stale or weak sensor value shall not be presented as a
valid measurement. Health-related output shall be labelled as coaching support,
not diagnosis.

\paragraph{Security Requirements}
Session commands and relay routes shall require a role or private token. Secrets
shall stay outside tracked firmware and frontend source. The relay shall check
the token, source, timestamp, value range, and body size, and shall create no
persistent storage. Personal
session records shall remain on the project PC unless the team makes an
authorised export.

\subsection{Cost Analysis}
Table~\ref{tab:cost_analysis} gives a planning estimate for five wearables. The
project PC, domain, and VPS are already available, so they add no new purchase
cost. Invoice values can replace these estimates in the final report.

\begin{table}[H]
\centering
\caption{Estimated cost for a five-player prototype}
\label{tab:cost_analysis}
\small
\begin{tabularx}{\textwidth}{X r r r}
\toprule
\textbf{Item} & \textbf{Qty.} & \textbf{Unit cost (Rs.)} & \textbf{Amount (Rs.)}\\
\midrule
ESP32-S3 development board & 5 & 900 & 4,500\\
MPU-6050 IMU module & 5 & 180 & 900\\
MAX30102 optical sensor module & 5 & 350 & 1,750\\
NEO-6M GPS module and antenna & 5 & 650 & 3,250\\
Battery, charger, and regulation & 5 & 550 & 2,750\\
Enclosure, strap, wiring, and connectors & 5 & 650 & 3,250\\
Camera support and tripod allowance & 1 & 3,000 & 3,000\\
Calibration and labelling material & 1 & 1,200 & 1,200\\
\midrule
\multicolumn{3}{r}{\textbf{Estimated total}} & \textbf{20,600}\\
\bottomrule
\end{tabularx}
\end{table}

One wearable has an estimated material cost of Rs. 3,280. The design can start
with fewer units for a controlled drill and later increase to five. Annotation,
testing, fitting, and maintenance still require team effort.

\subsection{Risk Analysis}
\begin{longtable}{@{}p{1.1in}>{\centering\arraybackslash}p{0.65in}>{\centering\arraybackslash}p{0.6in}p{3.35in}@{}}
\caption{Project risk register}\label{tab:risk_register}\\
\toprule
\textbf{Risk} & \textbf{Chance} & \textbf{Impact} & \textbf{Control and check}\\
\midrule
\endfirsthead
\toprule
\textbf{Risk} & \textbf{Chance} & \textbf{Impact} & \textbf{Control and check}\\
\midrule
\endhead
Player ID switch & Medium & High & Team calibration, jersey and appearance votes, unresolved state, manual correction, IDF1 and switch count.\\
Ball missed & High & Medium & Ball-specific labels, short recovery, speed and field limits, visible confidence, review clips.\\
PPG motion noise & High & High & Stable fitting, contact and motion gates, raw signal retention, reference-device test.\\
GPS dropout & Medium & Medium & Fix and satellite state, physical bounds, vision cross-check, no long silent interpolation.\\
Clock drift & Medium & High & SNTP source time, receive time, sequence numbers, and repeated LED plus tap test.\\
Limited outcome labels & High & High & Minimum-data gate, reviewed labels, later-session test, no early injury claim.\\
Relay outage & Medium & High & Bounded wearable queue, domain health check, sequence replay, fresh snapshot after a long gap.\\
Camera movement & Medium & High & Calibration lock, field-point check, mapping-error monitor, recalibration step.\\
Dashboard overload & Medium & Medium & Large pitch, short labels, task-based pages, projector and narrow-screen review.\\
Unauthorised access & Low & High & TLS, token checks, local records, no relay database, limited logs.\\
\bottomrule
\end{longtable}
